\chapterbackmatter{Solutions to Supplementary Exercises}

\begin{enumerate}
  \SupplementarySolution{1}{Goldstone Directions in the \(O(N)\) Linear Sigma Model}
  \begin{enumerate}
    \item[(a)]
    The transformations that preserve \(\boldsymbol e_N\) rotate the first
    \(N-1\) components among themselves, so
    \[
      H=O(N-1),\qquad
      \mathcal M_{\rm vac}=\frac{O(N)}{O(N-1)}\simeq S^{N-1}.
    \]
    Therefore
    \[
      \dim O(N)-\dim O(N-1)
        =\frac{N(N-1)}2-\frac{(N-1)(N-2)}2=N-1.
    \]
    The broken generators are \(T_{aN}\), \(a=1,\ldots,N-1\), which rotate
    the chosen vacuum direction into each independent tangent direction.
    There is one Goldstone field for each of them.
    \item[(b)]
    The displacement from the minimum is
    \[
      \boldsymbol\phi^2-v^2
        =2v\sigma+\sigma^2+\boldsymbol\pi^2.
    \]
    Thus the quadratic part of the potential is
    \[
      V^{(2)}=\lambda v^2\sigma^2
        =\frac12(2\lambda v^2)\sigma^2.
    \]
    Hence
    \[
      \boxed{m_\sigma^2=2\lambda v^2,\qquad
        m_{\pi_a}^2=0\quad(a=1,\ldots,N-1).}
    \]
    The massive \(\sigma\) is normal to the vacuum sphere, while the
    \(\pi_a\) point along its flat tangent directions.  These are exactly
    the \(N-1\) directions generated by \(O(N)/O(N-1)\).
    \item[(c)]
    Infinitesimal invariance says
    \[
      \frac{\partial V}{\partial\phi_n}(t\phi)_n=0.
    \]
    Differentiate with respect to \(\phi_k\) and then set \(\phi=v\):
    \[
      \left.\frac{\partial^2V}{\partial\phi_k\partial\phi_n}\right|_v
        (tv)_n
      +\left.\frac{\partial V}{\partial\phi_n}\right|_v t_{nk}=0.
    \]
    The second term vanishes because \(v\) is stationary.  With
    \((M^2)_{kn}=\partial_k\partial_nV|_v\),
    \[
      \boxed{(M^2)_{kn}(tv)_n=0.}
    \]
    If \(tv\neq0\), the generator is broken and supplies a nonzero null
    vector of the Hessian, hence a massless scalar mode.
  \end{enumerate}

  \SupplementarySolution{2}{Chiral Symmetry Breaking, Mass Spurions, and \(U(1)_A\)}
  \begin{enumerate}
    \item[(a)]
    Under \((L,R)\in SU(N_f)_L\times SU(N_f)_R\), the condensate matrix
    transforms as
    \[
      \Sigma\mathbf1\longmapsto L(\Sigma\mathbf1)R^\dagger.
    \]
    It is unchanged precisely when \(LR^\dagger=\mathbf1\), or \(L=R\).
    The unbroken group is therefore \(SU(N_f)_V\).  The number of broken
    generators is
    \[
      2(N_f^2-1)-(N_f^2-1)=N_f^2-1.
    \]
    Thus two light flavors give three Goldstones and three light flavors
    give eight.
    \item[(b)]
    The assigned transformations give
    \[
      M^\dagger U\longmapsto
        LM^\dagger R^\dagger RUL^\dagger
        =L(M^\dagger U)L^\dagger,
    \]
    and similarly for \(U^\dagger M\), so the trace is invariant.  Since
    \[
      U=1+\frac{i\tau_a\pi_a}{f}
        -\frac{\boldsymbol\pi^2}{2f^2}+\cdots,
    \]
    the parity-even combination is
    \[
      \operatorname{Tr}(M^\dagger U+U^\dagger M)
        =2(m_u+m_d)
         -\frac{m_u+m_d}{f^2}\boldsymbol\pi^2+\cdots.
    \]
    The term linear in \(\tau_a\pi_a\), which could distinguish \(m_u\)
    from \(m_d\), cancels between \(U\) and \(U^\dagger\).  Consequently all
    three pions receive the same first-order mass proportional to
    \(m_u+m_d\); isospin breaking starts at higher order for their masses.
    \item[(c)]
    The non-singlet breaking
    \(SU(3)_L\times SU(3)_R\to SU(3)_V\) gives eight modes.  If an additional
    exact \(U(1)_A\) were also spontaneously broken, its generator would
    provide one more, for a total of nine pseudoscalar Goldstones.  The
    observed light multiplet is only the octet.  The failed premise is the
    exactness of \(U(1)_A\): the quantum gauge theory has a singlet axial
    anomaly, so this is not a conserved symmetry to which Goldstone's
    theorem can be applied.
  \end{enumerate}

  \SupplementarySolution{3}{Nonlinear Sigma-Model Geometry and Soft Zeros}
  \begin{enumerate}
    \item[(a)]
    With the overall factor \(F^2\) appropriate to a sphere of radius \(F\),
    \[
      F^2\partial_\mu n\cdot\partial^\mu n
        =\partial_\mu\boldsymbol\pi\cdot\partial^\mu\boldsymbol\pi
         +\frac{(\boldsymbol\pi\cdot\partial_\mu\boldsymbol\pi)
                (\boldsymbol\pi\cdot\partial^\mu\boldsymbol\pi)}
               {F^2-\boldsymbol\pi^2}.
    \]
    Hence
    \[
      \boxed{g_{ab}(\pi)
        =\delta_{ab}+\frac{\pi_a\pi_b}{F^2-\boldsymbol\pi^2}
        =\delta_{ab}+\frac{\pi_a\pi_b}{F^2}+O(\pi^4/F^4).}
    \]
    The interactions are derivative interactions because this metric
    depends on position on the target sphere but there is no potential.
    \item[(b)]
    The Baker--Campbell--Hausdorff expansion gives
    \[
      \gamma^{-1}\partial_\mu\gamma
        =i t_aD_{a\mu}+it_3E_\mu,
    \]
    with
    \[
      D_{a\mu}
        =\partial_\mu\xi_a
         -\frac16\left[
          \boldsymbol\xi^2\partial_\mu\xi_a
          -\xi_a(\boldsymbol\xi\cdot\partial_\mu\boldsymbol\xi)
         \right]+O(\xi^4\partial\xi),
    \]
    \[
      E_\mu=\frac12
        (\xi_1\partial_\mu\xi_2-\xi_2\partial_\mu\xi_1)
        +O(\xi^3\partial\xi).
    \]
    Under a general \(SO(3)\) transformation, \(D_{a\mu}\) is rotated
    homogeneously by the compensating \(SO(2)\), whereas \(E_\mu\) acquires
    the derivative of that local compensator.  Thus \(D\) is a covariant
    matter-like object and \(E\) is a connection.
    \item[(c)]
    LSZ reduction replaces the external Goldstone by its momentum \(q\).
    Since every local vertex containing that leg supplies at least one
    derivative, a diagram with no singularity at \(q=0\) is proportional to
    \(q\) and vanishes:
    \[
      \mathcal M(\cdots,\pi^a(q))=O(q).
    \]
    The qualification is an emission from an external on-shell line.  The
    neighboring propagator can behave as
    \(1/(2p\cdot q)\), canceling the explicit soft momentum.  The full soft
    theorem therefore consists of the Adler zero plus separately evaluated
    external-leg pole terms.
  \end{enumerate}

  \SupplementarySolution{4}{Goldstone Power Counting, Curvature, and Pseudo-Goldstones}
  \begin{enumerate}
    \item[(a)]
    Direct counting gives
    \[
      \nu=4L-2I+\sum_iV_id_i.
    \]
    For a connected graph,
    \(L=I-\sum_iV_i+1\), so \(I=L+\sum_iV_i-1\).
    Substitution yields
    \[
      \boxed{\nu=2+2L+\sum_iV_i(d_i-2).}
    \]
    A one-loop graph with only \(d_i=2\) vertices has
    \(\nu=2+2=4\).  It competes with a tree graph containing one
    four-derivative vertex.
    \item[(b)]
    Substitution of the normal-coordinate expansion into the kinetic term
    gives
    \[
      \mathcal L_4
        =\frac16R_{acbd}\pi_c\pi_d
          \partial_\mu\pi_a\partial^\mu\pi_b.
    \]
    Thus the four-point vertex is entirely determined by \(R_{acbd}\).
    For a sphere of radius \(F\),
    \[
      R_{acbd}=\frac1{F^2}
        (\delta_{ab}\delta_{cd}-\delta_{ad}\delta_{cb}).
    \]
    On shell, the resulting crossing-symmetric amplitude is
    \[
      \boxed{\mathcal M_{ab,cd}
        =\frac1{F^2}\left(
         \delta_{ab}\delta_{cd}s+
         \delta_{ac}\delta_{bd}t+
         \delta_{ad}\delta_{bc}u\right).}
    \]
    \item[(c)]
    Choose the aligned vacuum \(n_3=1\) and parameterize
    \[
      n_3=\sqrt{1-\boldsymbol\pi^2/F^2}
        =1-\frac{\boldsymbol\pi^2}{2F^2}+\cdots.
    \]
    The perturbation becomes
    \[
      \Delta\mathcal L=hF-\frac{h}{2F}\boldsymbol\pi^2+\cdots.
    \]
    Comparing the quadratic term with
    \(-m^2\boldsymbol\pi^2/2\) gives
    \[
      \boxed{m_{\pi_1}^2=m_{\pi_2}^2=h/F.}
    \]
    The squared mass is linear in the explicit-breaking source, the general
    pseudo-Goldstone behavior.
  \end{enumerate}

  \SupplementarySolution{5}{Current Poles, PCAC, and Pion Thresholds}
  \begin{enumerate}
    \item[(a)]
    Lorentz covariance lets the part of the Fourier transform relevant to
    a scalar intermediate state be written
    \[
      G^\mu_{ab}(q)=q^\mu G_{ab}(q^2).
    \]
    Current conservation would imply \(q^2G_{ab}=0\) away from contact
    terms.  The equal-time contact term is the nonzero expectation value
    \(\langle[J^0_a,\Phi_b]\rangle\), so \(G_{ab}\) cannot vanish
    identically.  It must instead have
    \[
      G_{ab}(q^2)\sim\frac{C_{ab}}{q^2-i0}.
    \]
    The spectral density therefore contains a contribution proportional to
    \(\delta(q^2)\).  This is a massless one-particle state coupled both to
    the broken current and to \(\Phi_b\).
    \item[(b)]
    Translation invariance and the on-shell condition give, up to the
    overall phase convention for one-particle states,
    \[
      \boxed{
      \bra{\Omega}\partial_\mu A^\mu_a(0)\ket{\pi_b(p)}
        =f_\pi m_\pi^2\delta_{ab}.}
    \]
    Thus the divergence is small in precisely the pseudo-Goldstone
    parameter \(m_\pi^2\).  In the chiral limit \(m_\pi\to0\), the axial
    current is conserved while its one-pion matrix element remains
    nonzero.
    \item[(c)]
    At threshold \(s=4m_\pi^2\), \(t=u=0\), so
    \[
      A(s,t,u)=\frac{3m_\pi^2}{f_\pi^2},\qquad
      A(t,s,u)=A(u,t,s)=-\frac{m_\pi^2}{f_\pi^2}.
    \]
    Hence
    \[
      T^0_{\rm thr}=\frac{7m_\pi^2}{f_\pi^2},\qquad
      T^1_{\rm thr}=0,\qquad
      T^2_{\rm thr}=-\frac{2m_\pi^2}{f_\pi^2}.
    \]
    The \(S\)-wave scattering lengths are therefore
    \[
      \boxed{a_0^0=\frac{7m_\pi^2}{32\pi f_\pi^2},\qquad
        a_0^2=-\frac{m_\pi^2}{16\pi f_\pi^2}.}
    \]
    Bose symmetry excludes an \(I=1\) \(S\)-wave at threshold.
  \end{enumerate}

  \SupplementarySolution{6}{Goldberger--Treiman, Baryon Counting, and Soft-Pion Commutators}
  \begin{enumerate}
    \item[(a)]
    Let \(q=p'-p\).  Anticommuting \(\gamma_5\) through the Dirac matrices
    and using the two on-shell equations gives
    \[
      \overline u(p')\slashed q\gamma_5u(p)
        =2m_N\overline u(p')\gamma_5u(p),
    \]
    up to the conventional factors of \(i\) associated with the chapter's
    Dirac conventions.  The derivative vertex is therefore on-shell
    equivalent to
    \[
      -\frac{g_Am_N}{f_\pi}
        \overline N\gamma_5\tau_a\pi_aN.
    \]
    Identifying the coefficient with the pseudoscalar pion--nucleon
    coupling gives
    \[
      \boxed{g_{\pi NN}f_\pi=g_Am_N.}
    \]
    In the chapter's convention \(F_\pi=2f_\pi\), this is
    \(G_{\pi N}=2m_Ng_A/F_\pi\).
    \item[(b)]
    Before using topology,
    \[
      \nu=4L-2I_\pi-I_N+\sum_iV_id_i.
    \]
    The connected-graph identities are
    \[
      L=I_\pi+I_N-\sum_iV_i+1,\qquad
      2I_N+E_N=\sum_iV_in_i.
    \]
    Eliminating \(I_\pi\) and \(I_N\) gives
    \[
      \boxed{\nu
        =2+2L-E_N+\sum_iV_i
         \left(d_i+\frac{n_i}{2}-2\right).}
    \]
    Pure-pion vertices have \(n_i=0,d_i\geq2\), and leading
    pion--nucleon vertices have \(n_i=2,d_i\geq1\), so every bracket is
    nonnegative.
    \item[(c)]
    Contract the axial-current identity with \(q_\mu\), isolate its pion
    pole, and apply LSZ to the outgoing pion.  With
    \(\langle\Omega|A^\mu_a|\pi_b(q)\rangle=if_\pi q^\mu\delta_{ab}\),
    the regular part gives
    \[
      \boxed{
      \lim_{q\to0}
      \langle\beta;\pi^a(q)|\mathcal O|\alpha\rangle
        =-\frac{i}{f_\pi}
         \langle\beta|[Q_5^a,\mathcal O]|\alpha\rangle
         +\text{external-leg poles}.}
    \]
    Reversing the phase convention for the pion field reverses both sides,
    so physical consequences are unchanged.  Pole terms must be retained
    when the soft current can attach to an on-shell external particle.
  \end{enumerate}

  \SupplementarySolution{7}{Chiral Invariants, Mass Spurions, and Gell-Mann--Okubo}
  \begin{enumerate}
    \item[(a)]
    The product transforms by conjugation:
    \[
      \partial_\mu U\,\partial^\mu U^\dagger
        \longmapsto
        R(\partial_\mu U\,\partial^\mu U^\dagger)R^\dagger,
    \]
    so its trace is invariant.  A single derivative cannot form a
    parity-even Lorentz scalar, while two derivatives and unitarity reduce
    every single-trace candidate to this one; a two-trace term vanishes
    because \(\operatorname{Tr}(U^\dagger\partial_\mu U)=0\) for
    \(U\in SU(3)\).

    If \(U=\exp(2iB/F)\), then
    \[
      \operatorname{Tr}(\partial_\mu U\partial^\mu U^\dagger)
        =\frac4{F^2}\operatorname{Tr}
         (\partial_\mu B\partial^\mu B)+O(B^4/F^4).
    \]
    Choosing the overall coefficient accordingly produces canonical
    kinetic terms for the eight real fields in the traceless Hermitian
    matrix \(B\).
    \item[(b)]
    The invariant is
    \[
      \boxed{\operatorname{Tr}(\chi^\dagger U+U^\dagger\chi).}
    \]
    Each product transforms by conjugation under one chiral factor, and the
    displayed sum is even under parity.  At linear order in \(\chi\), any
    other ordering is equivalent by cyclicity of the trace.

    Setting \(\chi\) equal to the quark-mass matrix generates a
    pseudo-Goldstone mass term.  Since \(m_{\rm GB}^2\propto m_q\) while a
    derivative counts as \(Q\), consistency of the low-energy expansion
    requires
    \[
      \boxed{\chi\sim m_q\sim Q^2.}
    \]
    \item[(c)]
    Direct substitution gives
    \[
      4m_K^2-m_\pi^2
        =4B_0(\widehat m+m_s)-2B_0\widehat m
        =2B_0(\widehat m+2m_s).
    \]
    The right-hand side equals \(3m_\eta^2\).  Therefore
    \[
      \boxed{3m_\eta^2+m_\pi^2=4m_K^2.}
    \]
    The relation applies to squared masses because quark masses generate
    pseudo-Goldstone mass squared at leading chiral order.
  \end{enumerate}

  \SupplementarySolution{8}{Quark-Mass Ratios, Loop Counterterms, and the WZW Term}
  \begin{enumerate}
    \item[(a)]
    Write the common charged electromagnetic shift as \(\Delta\).
    The leading formulas are
    \[
    \begin{aligned}
      m_{\pi^+}^2&=B_0(m_u+m_d)+\Delta,&
      m_{\pi^0}^2&=B_0(m_u+m_d),\\
      m_{K^+}^2&=B_0(m_u+m_s)+\Delta,&
      m_{K^0}^2&=B_0(m_d+m_s).
    \end{aligned}
    \]
    Eliminating \(B_0,\Delta,m_s\) yields
    \[
      \boxed{
      \frac{m_u}{m_d}
        =\frac{
         2m_{\pi^0}^2-m_{K^0}^2-m_{\pi^+}^2+m_{K^+}^2}
        {m_{K^0}^2+m_{\pi^+}^2-m_{K^+}^2}.}
    \]
    Corrections to Dashen's leading electromagnetic relation and higher
    chiral orders modify the numerical value.
    \item[(b)]
    The pure-Goldstone counting formula is
    \[
      \nu=2+2L+\sum_iV_i(d_i-2).
    \]
    A one-loop graph with only \(d_i=2\) vertices has \(\nu=4\).  A tree
    graph with one \(d_i=4\) vertex also has
    \(\nu=2+(4-2)=4\).  They must therefore be included together.

    The loop contains ultraviolet poles and logarithms proportional to
    local \(Q^4\) operators.  Renormalized coefficients \(L_i^r(\mu)\) in
    \(\mathcal L_4\) absorb the poles and run so that their
    \(\ln\mu\) dependence cancels the loop logarithms.  Only the sum is a
    scale-independent amplitude.
    \item[(c)]
    At lowest order,
    \(U^{-1}dU=(2i/F)dB+\cdots\), so the five-form starts with
    \[
      \operatorname{Tr}(dB\,dB\,dB\,dB\,dB).
    \]
    Converting the five-ball integral to its four-dimensional boundary
    gives a term of the form
    \[
      \epsilon^{\mu\nu\rho\sigma}
        \operatorname{Tr}
        (B\,\partial_\mu B\,\partial_\nu B\,
         \partial_\rho B\,\partial_\sigma B).
    \]
    Thus the leading interaction contains five meson fields, four
    four-dimensional derivatives, and the parity-odd epsilon tensor.
    Two different extensions of the same boundary field differ by a map
    from \(S^5\) to \(SU(3)\).  Since \(\pi_5(SU(3))=\mathbb Z\), the action
    changes by \(2\pi\) times the coefficient; single-valuedness of
    \(e^{iI}\) requires that coefficient to be an integer.
  \end{enumerate}

  \SupplementarySolution{9}{Vacuum Alignment, CCWZ, and Residual Subgroups}
  \begin{enumerate}
    \item[(a)]
    Rotational invariance of the unperturbed potential makes the
    source-dependent energy smallest when
    \[
      \langle\boldsymbol\phi\rangle=\rho\,
        \frac{\boldsymbol h}{|\boldsymbol h|}.
    \]
    The radial stationarity equation is
    \(V_0'(\rho)=|\boldsymbol h|\).  A transverse fluctuation changes
    \(\boldsymbol\phi^2\) only at second order, and its curvature is
    \(V_0'(\rho)/\rho\).  Hence
    \[
      \boxed{m_\perp^2=\frac{|\boldsymbol h|}{\rho}
        =\frac{|\boldsymbol h|}{v}+O(h^2).}
    \]
    The source both selects a unique vacuum and turns the tangent modes
    into pseudo-Goldstones.
    \item[(b)]
    If \(\psi=\gamma(\xi)\widetilde\psi\) transforms linearly under \(G\),
    then
    \[
      g\psi=g\gamma(\xi)\widetilde\psi
        =\gamma(\xi')h(\xi,g)\widetilde\psi.
    \]
    Therefore
    \[
      \boxed{\widetilde\psi\longmapsto
        h(\xi,g)\widetilde\psi.}
    \]
    Decompose
    \(\gamma^{-1}d\gamma=iD^ax_a+iE^it_i\).
    Differentiating the factorization gives
    \[
      \boxed{iE'\equiv iE'^it_i
        =h(iE)h^{-1}-(dh)h^{-1}.}
    \]
    The inhomogeneous last term is precisely what cancels the derivative
    of \(h(\xi(x),g)\) in
    \((d+iE)\widetilde\psi\), so this covariant derivative transforms
    homogeneously.
    \item[(c)]
    The first-order vacuum energy must have the form
    \[
      E_1(v)=c\,v^\dagger Av+\text{constant}.
    \]
    Its extrema under \(v^\dagger v=1\) satisfy
    \(Av=\lambda v\), so alignment selects an extremal eigenvector of
    \(A\).  A generic Hermitian \(A\) with distinct eigenvalues has
    centralizer \(U(1)^{N-1}\) inside \(SU(N)\).  The spontaneous stabilizer
    additionally requires that the phase acting on the selected vector be
    one.  The determinant constraint then leaves
    \[
      \boxed{H_{\rm complete}=U(1)^{N-2}.}
    \]
    Degenerate eigenvalues enhance this residual group.
  \end{enumerate}

  \SupplementarySolution{10}{Positivity, \(U(1)_A\), and the Neutral-Pion Anomaly}
  \begin{enumerate}
    \item[(a)]
    The gauge covariant derivative is flavor blind, so \([\slashed D,M]=0\).
    Since \(\slashed D^\dagger=-\slashed D\),
    \[
      (\slashed D+M)^\dagger(\slashed D+M)
        =-\slashed D^{\,2}+M^2\geq m_{\min}^2.
    \]
    Every singular value of \(\slashed D+M\) is therefore at least
    \(m_{\min}\), and
    \[
      \boxed{\|(\slashed D+M)^{-1}\|\leq\frac1{m_{\min}}.}
    \]
    Correlators built from finitely many propagators are uniformly bounded
    as a small vector-symmetry-breaking perturbation is removed, provided
    no mass eigenvalue approaches zero.  There is then no singular factor
    capable of canceling that perturbation and leaving a symmetry-breaking
    order parameter.  The argument deliberately fails in a chiral
    zero-mass limit.
    \item[(b)]
    Classically,
    \[
      J^\mu_5=\sum_{f=1}^{N_f}\overline q_f\gamma^\mu\gamma_5q_f
    \]
    is conserved for massless quarks.  Quantum mechanically, with the
    conventional generator normalization,
    \[
      \boxed{\partial_\mu J^\mu_5
        =\frac{N_fg^2}{16\pi^2}
         F^a_{\mu\nu}\widetilde F^{a\mu\nu}.}
    \]
    Thus the singlet axial charge is not conserved.  Goldstone's theorem
    applies to the eight conserved non-singlet axial generators, but not to
    this anomalous ninth generator.  No extra light singlet pseudoscalar is
    required.
    \item[(c)]
    The vertex produces a gauge-invariant amplitude proportional to
    \[
      g_{\pi\gamma\gamma}
        \epsilon^{\mu\nu\rho\sigma}
        \epsilon_{1\mu}\epsilon_{2\nu}k_{1\rho}k_{2\sigma}.
    \]
    Summing the two photon polarizations and including the \(1/2!\) for
    identical particles gives
    \[
      \Gamma=\frac{g_{\pi\gamma\gamma}^2m_\pi^3}{64\pi}.
    \]
    With the anomalous coupling specified in the exercise,
    \[
      \boxed{\Gamma(\pi^0\to\gamma\gamma)
        =\frac{\alpha^2m_\pi^3}{64\pi^3f_\pi^2}.}
    \]
  \end{enumerate}

  \SupplementarySolution{11}{The Two-Dimensional Goldstone Obstruction}
  The Green function obeys
  \[
    -f^2\nabla^2G(x)=\delta^{(2)}(x).
  \]
  With ultraviolet reference length \(a\) and infrared scale \(L\),
  \[
    \boxed{G(x)=\langle X(x)X(0)\rangle
      =-\frac{1}{2\pi f^2}\log\frac{|x|}{L}}
  \]
  up to an additive constant.  The invariant field difference has variance
  \[
    \left\langle[X(x)-X(0)]^2\right\rangle
      =\frac{1}{\pi f^2}\log\frac{|x|}{a}.
  \]
  Thus a local displacement of a supposedly fixed vacuum does not die away:
  its fluctuations grow without bound at long distance.  Equivalently, the
  zero-momentum mode cannot be localized around one value of \(X\) in the
  infinite-volume limit.

  In \(D>2\) Euclidean dimensions the massless Green function falls as
  \(|x|^{2-D}\), so a local disturbance decays and an ordered vacuum is not
  ruled out by this infrared test.  In \(D=1\), corresponding to quantum
  mechanics, the Green function grows linearly with separation, an even
  stronger delocalization.  Two dimensions are the logarithmic borderline.

  If \(X\mapsto X+a\) were spontaneously broken, choosing a definite
  \(\langle X\rangle\) would produce a Goldstone field whose long-wavelength
  action is the one stated in the problem.  The logarithm shows that this
  assumed vacuum is washed out by its own infrared fluctuations.  The
  massless scalar theory may still exist as a two-dimensional conformal
  field theory, but its quanta are not Goldstone particles associated with
  distinct symmetry-breaking vacua.

\end{enumerate}
